\extrachap{Introduction to Part II}
\label{intro:pluripotential}
\addcontentsline{toc}{chapter}{Introduction to Part II}

In several complex variables, a positive Laplacian records too little information: it is only the trace of a Hessian whose positivity must hold in every complex direction. Complex differential forms retain that directional information and make it invariant under holomorphic coordinates. Currents extend the calculus to singular potentials. The first task of this part is to connect these invariant derivatives with the pointwise submean inequalities of Part I.

Once plurisubharmonic functions have been identified, the main difficulty is nonlinear. Measure coefficients cannot generally be multiplied. Local boundedness of the potential makes the recursive Bedford--Taylor construction possible, and decreasing continuity proves its symmetry and comparison properties. Its capacity detects pluripolar sets. The corresponding extremal functions, quasi-continuity and locality are developed before any attempt to remove the boundedness hypothesis, since all later truncation arguments depend on them.

There is a second route into singularities through holomorphic functions. Solving a weighted $\bar\partial$ equation corrects a cutoff error; singular weights force the correction to satisfy additional conditions. Extension from slices then controls integrability at a point and along a submanifold. These estimates lead to coherent multiplier ideals, Lelong level sets and strong openness. Weighted Bergman kernels turn the same information into approximations with analytic singularities. This order keeps the algebraic description of a singularity tied to the analytic estimates that justify it.

The passage to a compact K\"ahler manifold introduces a global constraint. A local potential can be changed by a pluriharmonic function on an overlap, and a prescribed cohomology class need not contain a smooth positive representative. Global potentials, positive classes and singular metrics provide the setting in which the local theory can still be used. Approximation and extension are established with their curvature losses and uniform constants before they enter the global constructions.

Non-pluripolar products are then formed by truncating the local potentials and retaining the product only where all factors are finite. Locality identifies the truncations on overlaps, while compact mass estimates make their limit finite. Removing the pole sets changes the mass problem: even a nonzero positive current can have zero first-degree non-pluripolar mass. Comparison, monotonicity and capacity convergence must therefore be proved for this new product, rather than inferred from bounded Stokes' formula.

Envelopes select potentials subject to obstacle and singularity constraints. Their contact properties lead to rooftop mass inequalities and model references. Relative capacities make integration by parts possible for potentials sharing an unbounded singularity, and energy then links the static envelope theory to equations and Perron segments. The ceiling identity in \cite[Theorem~1.2]{PSWZ26}, including its zero-mass case, is treated after the required relative equation theory. The final chapter gives the convex duality, real Monge--Amp\`ere measures and ample toric geometry needed to continue to the toric part of \cite{XiaSampleBook}.

The chapters therefore address two related questions throughout: what information about a singularity is retained by holomorphic integrability, and what information is retained by nonlinear mass? The local estimates and geometric constructions prepare both questions; the global products and envelopes distinguish their answers. Their later classification and applications are the subject of the continuation.
